\bisection{Experimental Design and Settings}{实验设计与设置}
\label{sec:experiments}
\bisubsection{Design principles}{设计原则}
\begin{paired}
\pair{The design separates primary family-level inference from secondary within-family analysis. The primary comparison evaluates all six DQN objectives against MPC-4 across the five fully coupled regimes, yielding 30 objective--regime contrasts under shared seed blocks and traces. Secondary comparisons examine how auxiliary supervision and Double-DQN target evaluation alter performance within the family. Standard DQN versus its centered full-action variant measures the effect of adding that auxiliary objective to an unchanged factual DQN backbone. Uncentered full-action Q and centered immediate-cost targets receive the same three modeled outcomes per state. They match the type and count of model outputs. However, they do not by themselves isolate centering or Bellman bootstrapping because their target scales differ and their auxiliary weights were not tuned separately. Double DQN with and without centered full-action supervision assesses whether the fixed auxiliary implementation has a similar descriptive effect when action selection and target evaluation are decoupled~\cite{vanhasselt2016double}. The $c\in\{0,0.5,1\}$ regimes and planners with $H\in\{1,2,4,6\}$ bound the family comparison as the implemented objective and planning depth change.}
{实验设计将主要的家族级推断与次级的家族内部分析分开。主要比较在五种完全耦合工况下评估全部六种 DQN 目标相对 MPC-4 的表现，在共享种子区组和轨迹下形成 30 个“目标—工况”对比。次级比较考察辅助监督和 Double-DQN 目标评估如何改变家族内部表现。标准 DQN 与其中心化全动作变体的比较用于衡量在事实 DQN 主干不变时加入该辅助目标的影响。非中心化全动作 Q 和中心化即时成本目标在每个状态下接收相同的三个模型化结果，模型输出的类型和数量也相同；但由于目标尺度不同且辅助权重未分别调参，它们本身不能单独识别中心化或 Bellman bootstrap 的作用。带或不带中心化全动作监督的 Double DQN 用于考察：当动作选择与目标评估解耦时，固定辅助实现是否具有相似的描述性影响~\cite{vanhasselt2016double}。$c\in\{0,0.5,1\}$ 的工况和 $H\in\{1,2,4,6\}$ 的规划器则界定了随着所实现目标与规划深度变化的家族比较边界。}
\pair{Standard DQN and its centered full-action variant share every component except the auxiliary objective. They use the same input, network, replay, exploration, optimizer, update cadence, target synchronization, and gradient clipping. The full-action Q control applies the same auxiliary weight directly to the three uncentered Bellman targets. The immediate-advantage control centers the negative analytical one-step costs but omits their bootstrapped future values. Both Double-DQN variants use online-network action selection with target-network evaluation in factual and auxiliary Bellman targets. Dueling networks, prioritized replay, multi-step returns, noisy networks, and distributional value learning remain outside the comparison.}
{标准 DQN 与其中心化全动作变体除辅助目标外共享全部组件，包括输入、网络、经验回放、探索策略、优化器、更新节奏、目标网络同步和梯度裁剪。全动作 Q 对照将相同辅助权重直接施加于三个未中心化的 Bellman 目标。即时优势对照对解析一步成本的相反数进行中心化，但不包含 bootstrap 的未来价值。两个 Double-DQN 变体在事实和辅助 Bellman 目标中均由在线网络选择动作、目标网络评估动作。Dueling 网络、优先经验回放、多步回报、噪声网络和分布式价值学习不在本次比较范围内。}
\pair{The auxiliary weight was fixed before the internally consistent simulator experiment. A preceding calibration with model seeds 700--703 compared $\lambda\in\{0.1,0.3,1,3\}$. It selected $\lambda=0.3$ by maximizing the worst relative improvement of the centered full-action variant across five coupled regimes, with the smaller weight as the tie-break. This selection rule was specific to the centered objective. The uncentered full-action, immediate-advantage, and Double-DQN auxiliary variants inherited $\lambda=0.3$ without method-specific calibration. Because auxiliary losses differ in scale and error structure, their comparisons are fixed-coefficient implementation comparisons rather than a factorial identification of centering or bootstrapping. The calibration used an earlier simulator version and is a prior specification rather than evidence from the current generator. Confirmation and all component comparisons use seeds 800--819 without post hoc removal, replacement, or selection. Every learned policy was trained only on the nominal $c=1$ distribution.}
{辅助权重在内部一致模拟器实验之前即被固定。先前使用模型种子 700--703 比较了 $\lambda\in\{0.1,0.3,1,3\}$，通过最大化中心化全动作变体在五种耦合工况下的最差相对改进来选择 $\lambda=0.3$，并在并列时优先选择较小权重。该规则仅针对中心化目标。非中心化全动作、即时优势和 Double-DQN 辅助变体未经各自校准，直接继承 $\lambda=0.3$。由于辅助损失的尺度和误差结构不同，这些比较是固定系数下的实现比较，而不是对中心化或 bootstrap 的析因识别。校准使用较早版本的模拟器，属于事先设定，不构成当前生成器的证据。确认实验及全部组件比较均使用种子 800--819，不进行事后删除、替换或筛选。所有学习策略仅在名义 $c=1$ 分布上训练。}
\end{paired}

\bisubsection{Training configuration}{训练配置}
\begin{paired}
\pair{Each network has 24 inputs, hidden widths 128--128--64, ReLU activations, and three outputs. Every model receives 600 episodes of 64 tasks, totaling 38,400 real interactions. Fixed settings are $\gamma=0.97$, learning rate $3\times10^{-4}$, batch size 128, replay capacity 100,000, and gradient-norm limit 10. Training begins after 128 transitions and performs one update every four interactions. The checkpoint field \texttt{target\_interval=250} is evaluated only on update calls; because updates occur at interaction counts divisible by four, synchronization occurs every 500 real interactions, equivalent to 125 optimizer updates. Epsilon decreases linearly from 1 to 0.05 over 13,440 interactions.}
{每个网络包含 24 个输入，隐藏层宽度为 128--128--64，采用 ReLU 激活，并输出三个动作值。每个模型训练 600 个回合，每回合 64 个任务，共计 38,400 次真实交互。固定设置为 $\gamma=0.97$、学习率 $3\times10^{-4}$、批量大小 128、经验回放容量 100,000、梯度范数上限 10。累积 128 个转移后开始训练，每四次交互更新一次。检查点字段 \texttt{target\_interval=250} 仅在更新调用时判断；由于更新发生在可被四整除的交互计数上，实际每 500 次真实交互同步一次，相当于 125 次优化器更新。Epsilon 在 13,440 次交互内从 1 线性下降至 0.05。}
\pair{The contextual bandit uses the same network, optimizer, exploration schedule, and interaction budget, with $\gamma=0$. It therefore learns the modeled one-step reward without a long-horizon Bellman term. Uniform replay stores only one-step factual transitions, with preview outputs attached as auxiliary training data. Standard DQN and Double DQN use no analytical alternative-action labels. The four auxiliary variants evaluate three modeled outcomes at every real interaction. The checkpoint manifest records these counts, optimizer updates, and wall time for every trained model.}
{上下文多臂老虎机使用相同的网络、优化器、探索计划和交互预算，但令 $\gamma=0$，因此只学习模型化的一步奖励，不含长时程 Bellman 项。均匀经验回放仅存储一步事实转移，预览输出作为辅助训练数据附加其中。标准 DQN 和 Double DQN 不使用解析备选动作标签；四个辅助变体在每次真实交互中评估三个模型化结果。检查点清单记录每个训练模型的这些计数、优化器更新次数和墙钟时间。}
\pair{All reported experiments use the internally consistent synthetic generator described above. Link bandwidth and contact follow bounded noisy periodic traces generated at reset. Burst and thermal-stress episodes multiply raw data and workload by 1.45 during one contiguous interval; all dependent heat, energy, latency, and transmission quantities are recomputed from those primitives. An earlier independently sampled generator is retained only as a historical reproducibility mode and is not used for the reported results.}
{所有报告实验均使用上述内部一致的合成生成器。链路带宽和过站时间遵循重置时生成的有界含噪周期轨迹。在突发和热应力回合中，原始数据量与工作负载在一个连续区间内乘以 1.45；全部相关热、能量、时延和传输量都由这些原语重新计算。较早的独立采样生成器仅作为历史复现模式保留，不用于本文结果。}
\end{paired}

\bisubsection{Seven evaluation regimes}{七种评估工况}
\begin{paired}
\pair{The seven regimes include two composite coupling-boundary conditions and five fully coupled operating conditions:
\begin{enumerate}[leftmargin=1.5em]
\item \textbf{Static ($c=0$):} removes action-dependent state transitions and scales delayed feasibility penalties to zero.
\item \textbf{Reduced coupling ($c=0.5$):} halves action-dependent resource effects and the delayed feasibility-penalty scale.
\item \textbf{Nominal ($c=1$):} matches the fully coupled training distribution.
\item \textbf{Burst:} increases contiguous compute and thermal demand.
\item \textbf{Link-limited:} reduces the available bandwidth trace by 0.22 before clipping.
\item \textbf{Energy-limited:} lowers initial energy from 0.88 to 0.62.
\item \textbf{Thermal-stress:} raises initial heat from 0.22 to 0.42 and includes the burst interval.
\end{enumerate}}
{七种工况包括两个复合耦合边界条件和五个完全耦合运行条件：
\begin{enumerate}[leftmargin=1.5em]
\item \textbf{静态（$c=0$）：}去除动作相关状态转移，并将延迟可行性惩罚缩放为零。
\item \textbf{弱耦合（$c=0.5$）：}将动作相关资源影响和延迟可行性惩罚尺度减半。
\item \textbf{名义（$c=1$）：}与完全耦合的训练分布一致。
\item \textbf{突发：}提高连续区间内的计算与热需求。
\item \textbf{链路受限：}在裁剪前将可用带宽轨迹降低 0.22。
\item \textbf{能量受限：}将初始能量从 0.88 降至 0.62。
\item \textbf{热应力：}将初始热状态从 0.22 提高至 0.42，并包含突发区间。
\end{enumerate}}
\pair{The first two regimes are composite interventions because the same $c$ appears in both the transition in Eq.~\eqref{eq:transition} and the delayed-penalty term in Eq.~\eqref{eq:reward}. They therefore do not isolate action-dependent carryover from reward rescaling. The remaining five regimes retain $c=1$. Models trained at $c=1$ are transferred to the two boundary conditions without retraining. These rows measure behavior under a joint objective-and-dynamics shift rather than the best policy attainable at lower coupling.}
{前两种工况属于复合干预，因为同一 $c$ 同时出现在式~\eqref{eq:transition} 的状态转移和式~\eqref{eq:reward} 的延迟惩罚项中。因此，它们不能将动作相关跨步延续效应与奖励缩放分离。其余五种工况均保持 $c=1$。在 $c=1$ 下训练的模型不经重新训练，直接迁移至两个边界条件。这两行衡量目标与动力学同时变化时的行为，而不是较低耦合下可达到的最优策略。}
\end{paired}

\bisubsection{Planning comparator, outcomes, and inference}{规划比较器、评价指标与统计推断}
\begin{paired}
\pair{The deployment comparison includes immediate-cost argmin, a contextual bandit, a prespecified resource-threshold heuristic, standard DQN, the centered full-action family member, and an exact $H$-step receding-horizon model-based planner. The threshold rule selects ground processing when $h>0.72$ or $e<0.24$ and both $b>0.28$ and $\tau>0.18$; under the same heat or energy condition it otherwise selects collaborative processing. If that condition is absent, it selects onboard processing when $b<0.22$ or $\tau<0.14$, collaborative processing when $q>0.70$, and immediate-cost argmin otherwise. These thresholds were hand specified and not fitted to the evaluation traces. Fixed onboard and fixed ground policies always select $\alpha_s$ and $\alpha_g$, respectively.}
{部署比较包括即时成本最小策略、上下文多臂老虎机、预设资源阈值启发式、标准 DQN、中心化全动作这一 DQN 家族成员，以及精确的 $H$ 步滚动时域模型规划器。阈值规则在 $h>0.72$ 或 $e<0.24$ 且同时满足 $b>0.28$ 与 $\tau>0.18$ 时选择地面处理；在相同热或能量条件下，若后两个条件不满足，则选择协同处理。若不存在上述热或能量条件，则在 $b<0.22$ 或 $\tau<0.14$ 时选择星上处理，在 $q>0.70$ 时选择协同处理，其他情况下选择即时成本最小动作。这些阈值由人工预设，未对评估轨迹拟合。固定星上和固定地面策略分别始终选择 $\alpha_s$ 和 $\alpha_g$。}
\pair{At every real decision, the planner enumerates all $3^H$ action sequences over the next $H$ known tasks and link states. It minimizes discounted modeled cost with $\gamma=0.97$, executes only the first action, and replans. The primary table retains $H=4$. A separate nominal analysis evaluates $H\in\{1,2,4,6\}$ on 400 unique held-out traces and reports wall-clock cost per decision. Each planner is exact only for its stated perfect-forecast objective, not for the 64-step global problem.}
{在每个真实决策时刻，规划器针对未来 $H$ 个已知任务和链路状态枚举全部 $3^H$ 个动作序列，以 $\gamma=0.97$ 最小化折扣模型成本，仅执行第一个动作，随后重新规划。主表保留 $H=4$。另一个名义分析在 400 条互不重复的留出轨迹上评估 $H\in\{1,2,4,6\}$，并报告每次决策的墙钟耗时。每个规划器仅对其所声明的完美预测有限时域目标是精确的，而不是 64 步全局问题的精确解。}
\pair{The primary outcome is the undiscounted episode total cost}
{主要评价指标为不折扣的回合总成本}
\end{paired}
\begin{equation}
C_{\mathrm{episode}}=\sum_{t=0}^{T-1}\left(c^{\mathrm{imm}}_t+p_t\right)=-\sum_{t=0}^{T-1}r_t,
\label{eq:episode_cost}
\end{equation}
\begin{paired}
\pair{where $p_t$ is the complete nonnegative penalty inside the square brackets in Eq.~\eqref{eq:reward}, including its multiplier $c$. Every executed step, including the terminal step, contributes once. DQN and the finite-horizon planner optimize discounted targets with $\gamma=0.97$, whereas Eq.~\eqref{eq:episode_cost} reports accumulated operational cost without discounting. This optimization--evaluation distinction was not subjected to a separate discount sensitivity analysis. Secondary outcomes are immediate cost, delayed penalty, energy use, modeled latency, transmitted data, threshold-exposure steps, planner wall-clock time, maximum heat, and minimum remaining energy.}
{其中，$p_t$ 是式~\eqref{eq:reward} 方括号内包含乘子 $c$ 的完整非负惩罚。每个已执行步骤（包括终止步）均计入一次。DQN 与有限时域规划器使用 $\gamma=0.97$ 优化折扣目标，而式~\eqref{eq:episode_cost} 报告不折扣的累积运行成本。该优化与评估目标的区别未单独进行折扣敏感性分析。次要指标包括即时成本、延迟惩罚、能量使用、模型时延、传输数据量、阈值暴露步数、规划器墙钟时间、最高热状态和最低剩余能量。}
\pair{The independent inferential unit is the model-seed block. For learned--learned comparisons, the same seed initializes network weights and controls episode task seeds, exploration draws, and replay-index sampling. Trajectories may diverge after different greedy actions, but the matched seed remains the randomization block. Each seed is evaluated on the same 20 held-out traces, which are averaged within seed, giving $n=20$ blocks per learned method. For the central DQN-family--MPC-4 claim, every raw row is matched on regime, model-seed trace namespace, and trace seed before trace-level differences are averaged within seed. MPC-4 is deterministic conditional on a trace; the model-seed field therefore indexes its matched trace namespace rather than a trained planner replicate.}
{独立推断单位是模型种子区组。对于学习方法之间的比较，相同种子用于初始化网络权重，并控制回合任务种子、探索抽样和回放索引采样。不同贪婪动作可能使轨迹分化，但匹配种子仍作为随机化区组。每个种子在相同的 20 条留出轨迹上评估，并先在种子内取平均，因此每种学习方法有 $n=20$ 个区组。对于核心的 DQN 家族与 MPC-4 比较，每一原始行先按工况、模型种子轨迹命名空间和轨迹种子匹配，再在种子内对轨迹级差值取平均。给定轨迹后 MPC-4 是确定性的，因此其模型种子字段索引的是匹配的轨迹命名空间，而非训练得到的规划器重复。}
\pair{All paired bootstrap analyses use 10,000 resamples with seed 20260829. Individual two-sided 95\% intervals are percentile intervals, and exact two-sided paired sign tests assess directional consistency. For the 30 DQN-objective--regime contrasts against MPC-4, one joint bootstrap resamples the same 20 seed blocks across all contrasts, preserving their correlation. Sign-test values are Holm-adjusted across all 30 cells. We additionally report the maximum, least-favorable mean difference within each regime and across all 30 cells. A simultaneous one-sided 95\% upper bound adds a joint critical value to every observed cell mean. The critical value is the 95th percentile of the largest centered bootstrap deviation across the relevant cells. An upper bound below zero supports the joint statement that every mean difference is negative.}
{全部配对 bootstrap 分析均使用随机种子 20260829 进行 10,000 次重采样。单个双侧 95\% 区间采用百分位区间，精确双侧配对符号检验评估方向一致性。对于 DQN 各目标相对 MPC-4 的 30 个“目标—工况”对比，一次联合 bootstrap 在所有对比中同步重采样相同的 20 个种子区组，从而保留其相关性。符号检验值在全部 30 个单元格上进行 Holm 校正。我们还报告每种工况内以及全部 30 个单元格中的最大、即最不利均值差。同时单侧 95\% 上界是在每个观测单元均值上加一个联合临界值；该临界值为相关单元格中最大中心化 bootstrap 偏差的第 95 百分位数。若上界低于零，则支持所有均值差均为负的联合陈述。}
\pair{For within-family contrasts, Holm adjustment is applied across the five fully coupled regimes separately for each prespecified contrast. The eight parameter profiles form a separate family, and each preview-scale contrast forms its own five-regime family. A directional cell is labelled supported only when its mean favors the first-named method, its paired interval excludes zero, and its family-adjusted sign-test value is below 0.05. We report seed-level differences and leave-one-seed-out ranges for the prespecified centered full-action versus standard-DQN contrast. Its conservative sensitivity analysis resamples the two sets of 20 seed means independently and ignores matching. Test traces and bootstrap draws are never treated as independent model replicates.}
{对于家族内部对比，每个预设对比分别在五种完全耦合工况上进行 Holm 校正。八个参数配置构成一个独立检验族，每个预览尺度对比各自构成包含五种工况的检验族。只有当均值有利于先列方法、配对区间不含零且家族校正后的符号检验值低于 0.05 时，方向性单元格才标记为“支持”。对于预设的中心化全动作与标准 DQN 对比，我们报告种子级差值和留一模型种子均值范围。其保守敏感性分析分别独立重采样两组各 20 个种子均值，不利用配对。测试轨迹和 bootstrap 抽样从不作为独立模型重复。}
\end{paired}

\bisubsection{Prespecified sensitivity and training diagnostics}{预设敏感性与训练诊断}
\begin{paired}
\pair{Without retraining or model selection, we repeated the nominal evaluation under eight profiles. They covered the reference configuration, compute demand $\pm20\%$, raw data volume $\pm20\%$, energy use $+20\%$, heat load $+20\%$, and link capacity $-20\%$. Each profile retained the same 20 models and the same held-out trace namespace as the primary experiment, with Holm adjustment across the eight profiles. Training logs record every episode for all 20 seeds. Curves report a causal 25-episode moving mean and seed-level mean $\pm$ SD for total cost, energy use, latency, and delayed penalty. These diagnostics describe optimization trajectories; confirmation claims remain based on held-out traces.}
{在不重新训练或选择模型的情况下，我们在八种配置下重复名义评估，包括参考配置、计算需求 $\pm20\%$、原始数据量 $\pm20\%$、能量使用 $+20\%$、热负载 $+20\%$ 和链路容量 $-20\%$。每种配置保留与主实验相同的 20 个模型及留出轨迹命名空间，并在八种配置上进行 Holm 校正。训练日志记录全部 20 个种子的每个回合。总成本、能量使用、时延和延迟惩罚曲线报告因果的 25 回合移动均值以及种子级均值 $\pm$ 标准差。这些诊断描述优化轨迹；确认性主张仍以留出轨迹为依据。}
\end{paired}

\bisubsection{Preview-model misspecification test}{预览模型误设检验}
\begin{paired}
\pair{The main experiment assumes that the analytical preview matches the simulator transition. To test this dependence, two additional centered full-action model sets were trained without changing the factual environment or deployment evaluation. One set scaled preview one-step costs (immediate cost plus modeled penalty) and action-induced resource changes by 0.85, while the other used 1.15. Resource predictions were clipped to the same physical bounds after scaling around the current resource state. Each condition used the same 20 seeds, training budget, network, replay, and held-out trace namespace as the main experiment. These scales were fixed as symmetric stress tests and were not used for model selection.}
{主实验假设解析预览与模拟器转移一致。为检验对该假设的依赖，在不改变事实环境或部署评估的情况下，另外训练两组中心化全动作模型。一组将预览的一步成本（即时成本加模型化惩罚）和动作引起的资源变化按 0.85 缩放，另一组按 1.15 缩放。以当前资源状态为中心完成缩放后，资源预测仍裁剪到相同物理边界。每个条件均使用与主实验相同的 20 个种子、训练预算、网络、经验回放和留出轨迹命名空间。这两个尺度作为对称压力检验预先固定，不用于模型选择。}
\end{paired}

\bisubsection{Reproducible result entry point}{可复现结果入口}
\begin{paired}
\pair{The core result entry point is \cmd{python -m dqn\_\allowbreak family\_\allowbreak satellite\_\allowbreak ground.\allowbreak reviewer\_\allowbreak experiments all --seeds 800-819 --resume}. Component and Double-DQN comparisons use \cmd{python -m dqn\_\allowbreak family\_\allowbreak satellite\_\allowbreak ground.\allowbreak extended\_\allowbreak revision --seeds 800-819 --stage all --resume}. The first command trains or loads 60 primary learned-policy checkpoints plus 40 preview-error checkpoints. The second reuses the primary DQN checkpoints and trains 80 component and Double-DQN checkpoints. It then evaluates six DQN objectives, performs the independent-seed bootstrap, and evaluates four planning depths on 400 unique traces.}
{核心结果入口为 \cmd{python -m dqn\_\allowbreak family\_\allowbreak satellite\_\allowbreak ground.\allowbreak reviewer\_\allowbreak experiments all --seeds 800-819 --resume}。组件与 Double-DQN 比较使用 \cmd{python -m dqn\_\allowbreak family\_\allowbreak satellite\_\allowbreak ground.\allowbreak extended\_\allowbreak revision --seeds 800-819 --stage all --resume}。第一条命令训练或加载 60 个主要学习策略检查点以及 40 个预览误差检查点。第二条命令复用主要 DQN 检查点，并训练 80 个组件与 Double-DQN 检查点；随后评估六种 DQN 目标，执行独立种子 bootstrap，并在 400 条互不重复的轨迹上评估四种规划深度。}
\pair{The command \cmd{python -m dqn\_\allowbreak family\_\allowbreak satellite\_\allowbreak ground.\allowbreak extended\_\allowbreak reporting} recomputes the 30 paired DQN--MPC-4 contrasts, joint maximum statistics, LaTeX tables, and artifact hashes from locked raw CSV files. The command \cmd{python -m dqn\_\allowbreak family\_\allowbreak satellite\_\allowbreak ground.\allowbreak transition\_\allowbreak audit} writes a fixed-seed, all-action example in which preview and executed costs and six resources agree coordinate by coordinate. Unit tests repeat this equality at both resource clipping boundaries and on the terminal step, and verify that time and link phase occupy the three added state coordinates.}
{命令 \cmd{python -m dqn\_\allowbreak family\_\allowbreak satellite\_\allowbreak ground.\allowbreak extended\_\allowbreak reporting} 从锁定的原始 CSV 文件重新计算 30 个配对 DQN--MPC-4 对比、联合最大统计量、LaTeX 表格和制品哈希。命令 \cmd{python -m dqn\_\allowbreak family\_\allowbreak satellite\_\allowbreak ground.\allowbreak transition\_\allowbreak audit} 写出一个固定种子的全动作示例，其中预览成本、执行成本及六个资源坐标逐项一致。单元测试在两个资源裁剪边界和终止步重复验证这一相等性，并确认时间与链路相位占据新增的三个状态坐标。}
\pair{Experiments used Python 3.13.9, NumPy 2.4.1, pandas 3.0.3, and PyTorch 2.8.0 with CUDA 12.9. Hardware comprised an NVIDIA GeForce RTX 5070 Ti running Windows 11. Planning wall time was measured on a 12th Gen Intel Core i9-12900K with 24 logical processors after one untimed warm-up episode. Values report the mean over all timed decisions. Machine-readable metadata record the 24-dimensional state definition, resolved variants, seeds, budgets, randomization blocks, model-information counts, result files, and runtime environment. Table and figure scripts consume these files directly.}
{实验使用 Python 3.13.9、NumPy 2.4.1、pandas 3.0.3，以及配套 CUDA 12.9 的 PyTorch 2.8.0。硬件为运行 Windows 11 的 NVIDIA GeForce RTX 5070 Ti。规划墙钟时间在配备 24 个逻辑处理器的第 12 代 Intel Core i9-12900K 上测量，并先运行一个不计时的预热回合；报告值为全部计时决策的均值。机器可读元数据记录 24 维状态定义、解析后的变体、种子、预算、随机化区组、模型信息计数、结果文件和运行环境。表格与图形脚本直接读取这些文件。}
\end{paired}
